\mode*
\section{The commit model}
\label{sec:model}

The first educational objective is Git's data model: commits as snapshots
of the whole tree, linked into a graph; history as that graph, not as a
sequence of diffs.

\subsection{Documented misconceptions and difficulties}

Learners' challenges with Git's conceptual model are documented from course
surveys and teaching observations \autocite{IsomottonenCochez2014}, and
Git's own concepts have been shown to misfit their purposes in analysable
ways \autocite{DeRossoJackson2016Gitless}.
% TODO(#1): The systematic phase should add repository-data studies of student
% Git usage and any misconception inventories beyond the seed sources
% (GitKit line of work, \autocite{Postner2026GitMisconceptions}).
% TODO(#2): Add misconceptions observed in our own course's grading data
% (datintro), clearly separated from literature-derived ones.

\subsection{Candidate critical aspects}

% XXX Preliminary; to be consolidated once the review phase is complete.
\begin{itemize}
  \item \emph{A commit is a snapshot, not a diff}: each commit contains the
    whole tree.
  \item \emph{History is a graph}: commits point at parents; nothing is a
    \enquote{folder of versions}.
  \item \emph{Nothing happens without a commit}: saving a file and
    recording a version are different acts.
\end{itemize}

\subsection{Patterns of variation}

\begin{itemize}
  \item Contrast (recorded vs not): edit a file, observe
    \texttt{git status}/\texttt{git log} before and after
    \texttt{git commit}.
  \item Generalisation (commit invariant, content varies): commits
    recording one file, many files, a deletion---always the same kind of
    object, shown by \texttt{git show}.
  \item Fusion: follow \texttt{git log --graph} while making a small
    history with a merge.
\end{itemize}
